\section{Proposta: arquitetura de segurança em camadas}
\label{sec:proposta}

Esta secção propõe uma arquitetura que responde às lacunas identificadas na
Secção~\ref{sec:resultados}. A proposta é apresentada ao nível do desenho e da sua
formalização; a avaliação experimental (custo de cada camada e ataques que cada uma
impede) é objeto da segunda parte do projeto. Uma implementação de referência em
C++23, com documentação técnica completa, está disponível no repositório do
projeto~\cite{mqttsecure2026repo}.

% ---------------------------------------------------------------------------
\subsection{Princípios de desenho}

\begin{enumerate}[label=\textbf{D\arabic*.}, leftmargin=*]
  \item \textbf{Defesa em profundidade}~(Eq.~\ref{eq:profundidade}): proteção no
    transporte, na mensagem e na disponibilidade, de forma independente.
  \item \textbf{O \emph{broker} não é de confiança para o conteúdo}: a
    confidencialidade, a integridade, a autenticidade e a frescura das mensagens
    não podem depender dele.
  \item \textbf{Camadas independentes e comutáveis}: cada mecanismo é uma camada
    que pode ser ativada ou desativada sem alterar as restantes, o que permite medir
    isoladamente o custo e o benefício de cada um.
  \item \textbf{Verificações baratas primeiro, estado só depois da autenticação}:
    rejeitar cedo o que é barato rejeitar e nunca deixar entrada não autenticada
    alterar estado.
  \item \textbf{Estado limitado}: toda a estrutura que um atacante consiga fazer
    crescer tem um limite superior.
\end{enumerate}

% ---------------------------------------------------------------------------
\subsection{Arquitetura em camadas}
\label{sec:arquitetura}

A comunicação de cada extremo (sensor, \emph{gateway}) passa por uma sequência
ordenada de camadas entre a aplicação e o transporte MQTT
(Figura~\ref{fig:camadas}). Seja $\mathcal{M}$ o conjunto das mensagens e
$\bot$ a rejeição. Cada camada $\ell_i$ define duas funções
$\ell_i^{\uparrow}, \ell_i^{\downarrow} : \mathcal{M} \to \mathcal{M} \cup \{\bot\}$,
com $\ell(\bot) = \bot$. Para camadas ordenadas do lado da aplicação ($\ell_1$)
para o lado do transporte ($\ell_n$):
\begin{align}
  \operatorname{Saída}(m)  &= \ell_n^{\downarrow} \circ \cdots \circ \ell_1^{\downarrow}(m), \label{eq:saida}\\
  \operatorname{Entrada}(m) &= \ell_1^{\uparrow} \circ \cdots \circ \ell_n^{\uparrow}(m). \label{eq:entrada}
\end{align}
A primeira rejeição interrompe o processamento: uma mensagem recebida que alguma
camada rejeite nunca chega à aplicação. A \emph{ordem} das camadas tem significado:
com a persistência acima da segurança guarda-se texto em claro; abaixo dela,
guardam-se envelopes (cifra em repouso); a limitação de tráfego fica junto ao
transporte para descartar inundações antes de qualquer operação criptográfica.

\begin{figure}[ht]
  \centering
  \begin{tikzpicture}[
      camada/.style={draw=iscte, fill=fundotab, rounded corners=2pt, minimum width=3.2cm,
                     minimum height=0.75cm, font=\small, align=center},
      ext/.style={draw=cinza, fill=white, rounded corners=2pt, minimum width=3.2cm,
                  minimum height=0.75cm, font=\small, align=center},
      seta/.style={-{Stealth[length=2mm]}, thick, iscte},
      nota/.style={font=\scriptsize, text=cinza, align=left, text width=7.4cm}]
    \node[ext] (app) {Aplicação};
    \node[camada, below=0.35cm of app] (pers) {Persistência (opcional)};
    \node[camada, below=0.35cm of pers] (sec) {\textbf{Segurança da mensagem}};
    \node[camada, below=0.6cm of sec] (lim) {Limitação de tráfego};
    \node[ext, below=0.35cm of lim] (mqtt) {Transporte MQTT\\\scriptsize TLS / mTLS};
    \node[ext, below=0.6cm of mqtt, fill=iscte!10] (brk) {\emph{Broker}\\\scriptsize autenticação, ACL, limites};
    \foreach \a/\b in {app/pers, pers/sec, sec/lim, lim/mqtt}
      \draw[seta] ([xshift=-0.6cm]\a.south) -- ([xshift=-0.6cm]\b.north);
    \foreach \a/\b in {pers/app, sec/pers, lim/sec, mqtt/lim}
      \draw[seta, acento] ([xshift=0.6cm]\a.north) -- ([xshift=0.6cm]\b.south);
    \draw[seta, <->] (mqtt) -- (brk);
    \node[font=\scriptsize, text=iscte, left=0.1cm of pers.west, anchor=east] {saída $\downarrow$};
    \node[font=\scriptsize, text=acento, right=0.1cm of pers.east, anchor=west] {$\uparrow$ entrada};
    \node[nota, right=0.9cm of sec] (n1) {Envelope fim-a-fim: AES-256-GCM, GMAC ou
      Ed25519; tópico nos dados associados; frescura; janela anti-repetição;
      política de modos; restrição chave--tópico.};
    \node[nota, right=0.9cm of lim] (n2) {\emph{Token bucket} por tópico,
      antes de qualquer operação criptográfica; estado limitado.};
    \node[nota, right=0.9cm of brk] (n3) {Camada 1: TLS 1.3, mTLS com identidade no
      certificado, palavras-passe, ACL por utilizador, identificador de cliente
      ligado ao utilizador, limites de pacote e de ligações.};
    \draw[cinza, densely dotted] (sec.east) -- (n1.west);
    \draw[cinza, densely dotted] (lim.east) -- (n2.west);
    \draw[cinza, densely dotted] (brk.east) -- (n3.west);
  \end{tikzpicture}
  \caption{Arquitetura proposta: camadas entre a aplicação e o transporte, cada uma
    comutável de forma independente. A saída percorre as camadas de cima para
    baixo e a entrada de baixo para cima.}
  \label{fig:camadas}
\end{figure}

% ---------------------------------------------------------------------------
\subsection{Camada de transporte e \emph{broker}}

A primeira camada usa o que os \emph{brokers} já oferecem
(Tabela~\ref{tab:seguranca-nativa}), configurado de forma restritiva: só a porta
TLS aberta (TLS~1.3); autenticação mútua por certificado, com o nome comum do
certificado como identidade MQTT; palavras-passe; ACL de negação por omissão com
regras parametrizadas pelo utilizador autenticado (por exemplo,
\texttt{restricted/sensors/\%u/\#} só para escrita); identificador de cliente
igual ao nome de utilizador, o que impede um cliente de tomar a sessão de outro;
desativação de \texttt{\$SYS}; e limites de tamanho de pacote e de número de
ligações. Esta camada protege a rede e restringe quem fala com o \emph{broker},
mas, por si só, deixa intactos os problemas P1 a P3.

% ---------------------------------------------------------------------------
\subsection{Envelope fim-a-fim}
\label{sec:envelope}

A camada de segurança substitui o conteúdo de cada mensagem por um
\emph{envelope} (Figura~\ref{fig:envelope}), com cabeçalho fixo de 25 octetos.

\begin{figure}[ht]
  \centering
  \begin{tikzpicture}[
      campo/.style={draw=iscte, minimum height=0.85cm, font=\scriptsize, align=center,
                    fill=fundotab},
      corpo/.style={draw=iscte, minimum height=0.85cm, font=\scriptsize, align=center,
                    fill=acento!20}]
    \node[campo, minimum width=0.9cm] (m)  {\texttt{"MS"}\\2};
    \node[campo, minimum width=0.8cm, right=0pt of m] (v) {ver.\\1};
    \node[campo, minimum width=0.9cm, right=0pt of v] (md) {modo\\1};
    \node[campo, minimum width=1.4cm, right=0pt of md] (k) {id. chave\\4};
    \node[campo, minimum width=1.8cm, right=0pt of k] (s) {sequência\\8};
    \node[campo, minimum width=1.8cm, right=0pt of s] (t) {instante (ms)\\8};
    \node[campo, minimum width=0.7cm, right=0pt of t] (n) {$|N|$\\1};
    \node[campo, minimum width=1.3cm, right=0pt of n] (nn) {\emph{nonce}\\$|N|$};
    \node[corpo, minimum width=3.6cm, right=0pt of nn] (b) {corpo\\variável};
    \draw[decorate, decoration={brace, amplitude=5pt}, thick, bom]
      ([yshift=2pt]m.north west) -- ([yshift=2pt]nn.north east)
      node[midway, above=6pt, font=\scriptsize, text=bom]
      {autenticado como dados associados, juntamente com o tópico MQTT};
    \draw[decorate, decoration={brace, amplitude=5pt, mirror}, thick, iscte]
      ([yshift=-2pt]m.south west) -- ([yshift=-2pt]n.south east)
      node[midway, below=6pt, font=\scriptsize, text=iscte] {cabeçalho fixo: 25 octetos};
  \end{tikzpicture}
  \caption{Formato do envelope (inteiros em \emph{big-endian}); o número sob cada
    campo é o seu tamanho em octetos.}
  \label{fig:envelope}
\end{figure}

\paragraph{Dados associados e ligação ao tópico.} Sendo $\mathit{cab}$ o cabeçalho,
$N$ o \emph{nonce} e $\tau$ o tópico MQTT, os dados associados são
\begin{equation}
  A = \mathit{cab} \,\|\, N \,\|\, \operatorname{u16}(|\tau|) \,\|\, \tau .
  \label{eq:aad}
\end{equation}
Como o tópico entra na autenticação mas não no envelope, uma mensagem válida
republicada noutro tópico falha a verificação sem qualquer custo de transmissão; e
como todo o cabeçalho (incluindo o modo, o identificador da chave, a sequência e o
instante) é autenticado, nenhum desses campos pode ser alterado.

\paragraph{Modos de proteção.} O remetente escolhe um de três modos, indicado no
cabeçalho:
\begin{align}
  \textbf{cifra:}\quad & \mathit{corpo} = C \,\|\, T, & (C, T) &= \operatorname{AES\text{-}GCM.Enc}(K, N, P, A); \label{eq:modo-cifra}\\
  \textbf{autenticação:}\quad & \mathit{corpo} = P \,\|\, T, & T &= \operatorname{AES\text{-}GCM.Enc}(K, N, \varnothing, A \,\|\, P); \label{eq:modo-auth}\\
  \textbf{assinatura:}\quad & \mathit{corpo} = P \,\|\, \sigma, & \sigma &= \operatorname{Ed25519.Sign}(\mathit{sk}, A \,\|\, P),\ N = \varnothing. \label{eq:modo-sign}
\end{align}
O modo de cifra garante confidencialidade e integridade; o de autenticação (GMAC)
garante integridade com o conteúdo legível; o de assinatura garante integridade e
\emph{prova de origem}: o recetor guarda apenas chaves públicas e, por isso, consegue
verificar, mas não forjar, mensagens dos dispositivos. O custo fixo por mensagem é
\begin{equation}
  \omega(\text{modo}) = 25 + |N| + |T| =
  \begin{cases}
    25 + 12 + 16 = 53 \text{ octetos} & \text{(cifra, autenticação)},\\
    25 + 0 + 64 = 89 \text{ octetos} & \text{(assinatura)}.
  \end{cases}
  \label{eq:overhead}
\end{equation}

\paragraph{Política do recetor contra a redução do nível de proteção.} Cada recetor
mantém o conjunto $\mathcal{A} \subseteq \{\text{cifra}, \text{autenticação},
\text{assinatura}\}$ dos modos que aceita e rejeita qualquer envelope cujo modo
$\mu \notin \mathcal{A}$. Como $\mu$ faz parte de $A$, um atacante não consegue
rotular uma mensagem cifrada como \enquote{apenas autenticada} sem invalidar a
etiqueta, e um detentor de chave não consegue baixar o nível de proteção abaixo do
que o recetor exige.

% ---------------------------------------------------------------------------
\subsection{Ordem de verificação no recetor}
\label{sec:ordem}

A Tabela~\ref{tab:ordem} define a sequência de verificações de um envelope
recebido. A ordem decorre do princípio D4.

\begin{table}[ht]
  \centering
  \caption{Sequência de verificações de um envelope recebido.}
  \label{tab:ordem}
  \small
  \begin{tabularx}{\textwidth}{@{}cL{3.5cm}YL{3.2cm}@{}}
    \toprule
    \textbf{\#} & \textbf{Verificação} & \textbf{Justificação da posição} & \textbf{Rejeição} \\
    \midrule
    1 & Formato do envelope & Descarta conteúdo não protegido (por exemplo, injeções em claro) & mensagem malformada \\
    2 & Modo $\mu \in \mathcal{A}$ & Barato; o modo está em $A$, pelo que mentir sobre ele só pode falhar & modo não aceite \\
    3 & Frescura, Eq.~\eqref{eq:frescura} & Barato; o instante está em $A$, pelo que alterá-lo invalida a etiqueta & mensagem expirada \\
    4 & Chave conhecida e verificação criptográfica & A partir daqui a mensagem é autêntica; a falha é sempre a mesma, sem revelar a causa (sem oráculo) & falha de autenticação \\
    5 & Restrição chave--tópico, Eq.~\eqref{eq:ligacao} & Só depois da autenticação: entrada não autenticada não pode sondar a política & permissão negada \\
    6 & Janela anti-repetição, Eq.~\eqref{eq:janela} & Só depois da autenticação (números forjados não avançam a janela) e da restrição (uma mensagem recusada não consome o seu número) & repetição detetada \\
    \bottomrule
  \end{tabularx}
\end{table}

% ---------------------------------------------------------------------------
\subsection{Frescura e anti-repetição}

A frescura usa a Eq.~\eqref{eq:frescura} com $\Delta_a = 60$~s e
$\Delta_s = 5$~s, e a anti-repetição a janela da Eq.~\eqref{eq:janela} com
$w = 64$, mantida por identificador de chave. Os dois mecanismos são
complementares: a frescura sozinha deixa passar todas as repetições com menos de
$\Delta_a$; a janela sozinha não protege contra a perda do seu estado (por exemplo,
reinício do recetor), situação em que a frescura limita a exposição a $\Delta_a$.
Para cumprir o princípio D5, o número de janelas é limitado e um identificador de
chave novo é recusado quando o limite é atingido. Os limites do GCM
(Eq.~\eqref{eq:aniversario}) e do espaço de sequências ($2^{64}$) impõem a rotação
das chaves.

% ---------------------------------------------------------------------------
\subsection{Restrição chave--tópico}
\label{sec:ligacao}

Autenticar uma mensagem prova apenas que foi produzida com a chave $k$. Para que um
dispositivo comprometido não publique, com a sua chave válida, nos tópicos dos
outros (problema P3), cada chave tem um dono $\operatorname{dono}(k)$ e o recetor
aplica regras de negação por omissão. Sendo $F_k$ as regras específicas da chave,
$F_\ast$ os padrões gerais e $\sigma_k$ a substituição
$\{\texttt{\%u} \mapsto \operatorname{dono}(k),\ \texttt{\%k} \mapsto k\}$:
\begin{equation}
  \operatorname{autorizada}(k, \tau) \iff
  \exists f \in F_k \cup \sigma_k(F_\ast) :\ \operatorname{corr}(f, \tau),
  \label{eq:ligacao}
\end{equation}
com $\operatorname{corr}$ definida na Eq.~\eqref{eq:corr}. Por exemplo, com o padrão
\texttt{+/sensors/\%u/\#}, a chave de \texttt{sensor-7} só pode publicar em
\texttt{open/sensors/sensor-7/\dots} e \texttt{restricted/sensors/sensor-7/\dots}.
O dono tem de ser um único nível de tópico sem carateres especiais, para que a
substituição não alargue o filtro.

% ---------------------------------------------------------------------------
\subsection{Gestão de chaves}

Cada dispositivo $d$ tem, por época $e$, uma chave AES, um par de chaves Ed25519 e o
identificador $\operatorname{id}(d, e) = e \cdot 10^{6} + d$ (o \emph{gateway} é
$d = 0$). Identificadores nunca são reutilizados: um identificador reutilizado teria
uma janela anti-repetição nova, na qual mensagens antigas voltariam a ser aceites. A
Tabela~\ref{tab:chaves} mostra que chaves cada papel detém.

\begin{table}[ht]
  \centering
  \caption{Distribuição de chaves por papel.}
  \label{tab:chaves}
  \small
  \begin{tabularx}{\textwidth}{@{}L{2.4cm}YYY@{}}
    \toprule
    \textbf{Papel} & \textbf{Chaves AES} & \textbf{Chaves privadas Ed25519} & \textbf{Chaves públicas Ed25519} \\
    \midrule
    \emph{Gateway}  & de todos os dispositivos e a sua & \textbf{apenas a sua} & de todos os dispositivos e a sua \\
    Dispositivo $d$ & a sua e a do \emph{gateway} (ordens) & a sua & a do \emph{gateway} \\
    Atacante        & --- & --- & --- (apenas o certificado da CA) \\
    \bottomrule
  \end{tabularx}
\end{table}

A \textbf{rotação} instala as chaves da época $e+1$ ao lado das da época $e$,
muda os remetentes para as novas e revoga as antigas após um período de graça. A
\textbf{revogação} de um dispositivo remove as suas chaves de mensagem, a sua
palavra-passe no \emph{broker} e coloca o seu certificado na CRL, atuando nas
duas camadas de defesa.

% ---------------------------------------------------------------------------
\subsection{Limitação de tráfego antes da criptografia}

Junto ao transporte, uma camada aplica a Eq.~\eqref{eq:token-bucket} por tópico,
rejeitando o excesso antes de qualquer verificação criptográfica. O número de baldes
é limitado; quando o limite é atingido, os baldes inativos (que voltariam a estar
cheios) são descartados e, se não houver nenhum, os tópicos novos partilham um único
balde, pelo que um atacante não consegue criar orçamento inventando tópicos.

% ---------------------------------------------------------------------------
\subsection{Mapeamento entre ameaças e mecanismos}

A Tabela~\ref{tab:ameacas-mecanismos} relaciona as ameaças da
Tabela~\ref{tab:stride} com os mecanismos da proposta.

\begin{table}[ht]
  \centering
  \caption{Ameaças e mecanismos da arquitetura proposta.}
  \label{tab:ameacas-mecanismos}
  \small
  \begin{tabularx}{\textwidth}{@{}L{5.0cm}YL{2.6cm}@{}}
    \toprule
    \textbf{Ameaça} & \textbf{Mecanismo} & \textbf{Camada} \\
    \midrule
    Escuta na rede               & TLS 1.3 & transporte \\
    Escuta pelo \emph{broker} ou por subscritores de \texttt{\#} & Modo de cifra; ACL & mensagem; \emph{broker} \\
    Adulteração do conteúdo ou do cabeçalho & Etiqueta GCM ou assinatura sobre $A \,\|\, P$ & mensagem \\
    Republicação noutro tópico   & Tópico em $A$, Eq.~\eqref{eq:aad} & mensagem \\
    Repetição                    & Janela anti-repetição & mensagem \\
    Mensagens antigas            & Frescura & mensagem \\
    Injeção de mensagens em claro & Formato do envelope e autenticação & mensagem \\
    Redução do nível de proteção & Modo em $A$ e política $\mathcal{A}$ & mensagem \\
    Usurpação de identidade no \emph{broker} & Palavras-passe, mTLS, identificador ligado ao utilizador & transporte \\
    Dispositivo comprometido a publicar por outros & Restrição chave--tópico, Eq.~\eqref{eq:ligacao}; ACL com \texttt{\%u} & mensagem; \emph{broker} \\
    Recetor a forjar dados de dispositivos & Modo de assinatura (recetor só com chaves públicas) & mensagem \\
    Chaves roubadas              & Rotação por épocas; revogação & gestão \\
    Inundação de um tópico       & \emph{Token bucket} antes da criptografia; limites no \emph{broker} & disponibilidade \\
    Esgotamento de memória       & Estado limitado (janelas, baldes, repositório) & todas \\
    Oráculo de verificação       & Falha uniforme; estado só depois da autenticação & mensagem \\
    \bottomrule
  \end{tabularx}
\end{table}

% ---------------------------------------------------------------------------
\subsection{Comparação com os trabalhos relacionados}
\label{sec:comparacao-proposta}

A Tabela~\ref{tab:proposta-vs} compara a proposta com os mecanismos fim-a-fim e de
autenticação mais próximos, acrescentando às propriedades da
Tabela~\ref{tab:comparacao} três mecanismos específicos desta proposta: ligação ao
tópico (LT), política contra a redução do nível de proteção (RN) e restrição
chave--tópico (CT).

\begin{table}[ht]
  \centering
  \caption{Posicionamento da proposta face aos trabalhos relacionados mais próximos.
    Siglas e símbolos como na Tabela~\ref{tab:comparacao}; LT ligação ao tópico;
    RN política contra a redução do nível de proteção; CT restrição chave--tópico;
    Com.: camadas comutáveis em tempo de execução.}
  \label{tab:proposta-vs}
  \scriptsize
  \setlength{\tabcolsep}{2.8pt}
  \renewcommand{\arraystretch}{1.2}
  \begin{tabular}{@{}lL{2.8cm}ccccccccc|cccc@{}}
    \toprule
    \textbf{ID} & \textbf{Trabalho}
      & \textbf{C} & \textbf{I} & \textbf{Au} & \textbf{Az} & \textbf{F/R} & \textbf{NR}
      & \textbf{D} & \textbf{P} & \textbf{BnC} & \textbf{LT} & \textbf{RN} & \textbf{CT} & \textbf{Com.} \\
    \midrule
    \E{02} & Singh et al.~\cite{singh2015secure}
      & \Sim & \Nr & \Nr & \Sim & \Nr & \Nr & \Nr & \Nr & \Parc & \Nr & \Nr & \Nr & \Nr \\
    \E{11} & Calabretta et al.~\cite{calabretta2018mqttauth}
      & \Sim & \Nr & \Sim & \Sim & \Nr & \Nr & \Nr & \Nr & \Nr & \Nr & \Nr & \Nr & \Nr \\
    \E{13} & Dinculeană e Cheng~\cite{dinculeana2019vulnerabilities}
      & \Sim & \Sim & \Nr & \Nr & \Nr & \Nr & \Nr & \Nr & \Sim & \Nr & \Nr & \Nr & \Nr \\
    \E{15} & Malina et al.~\cite{malina2019secure}
      & \Parc & \Sim & \Sim & \Nr & \Nr & \Nr & \Nr & \Sim & \Nr & \Nr & \Nr & \Nr & \Nr \\
    \E{16} & Sadio et al.~\cite{sadio2019lightweight}
      & \Sim & \Sim & \Nr & \Nr & \Nr & \Nr & \Nr & \Nr & \Sim & \Nr & \Nr & \Nr & \Nr \\
    \E{19} & Segarra et al.~\cite{segarra2020mqttz}
      & \Sim & \Nr & \Sim & \Nr & \Nr & \Nr & \Nr & \Nr & \Parc & \Nr & \Nr & \Nr & \Nr \\
    \E{20} & Dahlmanns et al.~\cite{dahlmanns2021entrust}
      & \Sim & \Sim & \Sim & \Nr & \Nr & \Nr & \Nr & \Nr & \Sim & \Nr & \Nr & \Nr & \Nr \\
    \E{26} & Javadi et al.~\cite{javadi2026cmap}
      & \Sim & \Sim & \Sim & \Nr & \Nr & \Nr & \Nr & \Nr & \Nr & \Nr & \Nr & \Nr & \Nr \\
    \midrule
    \rowcolor{fundotab}
    --- & \textbf{Proposta}
      & \Sim & \Sim & \Sim & \Sim & \Sim & \Sim & \Parc & \Nao & \Sim & \Sim & \Sim & \Sim & \Sim \\
    \bottomrule
  \end{tabular}
\end{table}

A proposta não supera os trabalhos relacionados em todas as dimensões: não oferece
anonimato~\cite{buccafurri2023mqtta}, não protege o próprio \emph{broker} com um
TEE~\cite{segarra2020mqttz}, não oferece confidencialidade de grupo por atributos
\cite{singh2015secure} e não tem verificação formal~\cite{javadi2026cmap}. O seu
contributo está na \emph{combinação}, numa só camada fim-a-fim, de propriedades
que a literatura analisada trata separadamente ou não reporta (ligação ao tópico,
frescura com anti-repetição, política contra a redução do nível de proteção,
restrição chave--tópico e não repúdio opcional) e na organização em camadas
comutáveis, que permite medir o custo e o benefício de cada mecanismo isoladamente.

% ---------------------------------------------------------------------------
\subsection{Limitações conhecidas}

\begin{itemize}[leftmargin=*]
  \item \textbf{Chaves pré-partilhadas}: a distribuição inicial das chaves assume um
    canal seguro; um acordo de chaves por dispositivo (por exemplo, ECDH) eliminaria
    os segredos pré-partilhados.
  \item \textbf{Chave de grupo nas ordens}: nos modos AES todos os atuadores detêm a
    chave do \emph{gateway}, pelo que um atuador comprometido pode forjar ordens para
    os outros; o modo de assinatura resolve-o, a um custo computacional superior.
  \item \textbf{Limitação por tópico}: o tópico é escolhido por quem publica, pelo
    que um atacante que inunde o tópico de uma vítima consome também o orçamento da
    vítima; limites por cliente pertencem ao \emph{broker}.
  \item \textbf{Estado da janela não persistente}: o reinício do recetor apaga as
    janelas anti-repetição, ficando a proteção limitada à frescura ($\Delta_a$).
  \item \textbf{Dados falsos mas autênticos}: um dispositivo comprometido continua a
    poder publicar dados falsos nos seus próprios tópicos; só a deteção e a
    revogação o travam.
\end{itemize}

% ---------------------------------------------------------------------------
\subsection{Plano de avaliação (segunda parte do projeto)}

A avaliação experimental medirá, para cada combinação de camadas: (i)~o custo, em
latência fim-a-fim, tempo de processador por mensagem, débito e octetos
acrescentados (Eq.~\ref{eq:overhead}); e (ii)~o benefício, como a fração de
ataques bem-sucedidos (escuta, injeção, repetição, mudança de tópico, usurpação de
identidade, inundação) com cada defesa ativa ou inativa. Esta medição conjunta de
custo e benefício responde diretamente à lacuna identificada na QF4.
